% CVPR 2026 submission draft — ReSPect
% Compiles with stock LaTeX (article class, CVPR-like layout).
% For submission, move into the official CVPR author kit.
\documentclass[10pt,twocolumn,letterpaper]{article}
\usepackage[margin=0.75in]{geometry}
\usepackage{graphicx,amsmath,amssymb,amsthm,booktabs}
\usepackage{hyperref}
\usepackage[numbers,sort&compress]{natbib}
\usepackage{caption}
\usepackage{xcolor}

\newtheorem{theorem}{Theorem}
\newtheorem{lemma}{Lemma}
\newtheorem{remark}{Remark}

\newcommand{\ie}{\textit{i.e.}}
\newcommand{\eg}{\textit{e.g.}}
\newcommand{\wrt}{\textit{w.r.t.}}

\title{ReSPect: Residual Spectral Prediction\\for Training-Free Acceleration of Diffusion Transformers}
\author{Anonymous CVPR Submission}
\date{}

\begin{document}
\maketitle

\begin{abstract}
Feature caching accelerates diffusion models by skipping denoiser
evaluations and filling in the missing features from history. The
literature has concentrated on \emph{when} to skip, culminating in
spectral-evolution-aware schedules that gate on Wiener-filtered content
drift. The complementary question---\emph{how} to reconstruct skipped
features---is still answered almost everywhere by zero-order reuse,
although reuse is, in a precise sense, the weakest predictor one can put
on the denoising trajectory. We introduce \textbf{ReSPect} (Residual
Spectral Prediction for Efficient Caching of Transformers), which keeps
the spectral gate but reconstructs skipped steps by forecasting
transformer \emph{residuals} with a ridge-regularized Chebyshev fit
blended against a discrete Taylor correction. Working in residual space
is not incidental: because the block input is known exactly on a skipped
step, all prediction error concentrates on the small centered residual,
which tightens both the approximation constant and the conditioning of
the fit. A bias--variance calculation selects the blend weight, and a
degrees-of-freedom argument sets the warm-up at $M{+}1$ observations,
below which the spectral fit is underdetermined and low-order prediction
is provably preferable. The resulting forecaster admits an error bound
with no dependence on skip length, in contrast to order-$P$ Taylor
extrapolation whose remainder grows as $|t_j - t_k|^{P+1}$. On
FLUX.1-dev (DrawBench, $1024{\times}1024$), ReSPect reaches
\textbf{27.97\,dB} PSNR against 26.29\,dB for the strongest reuse-based
schedule at matched compute, with matching gains in SSIM/LPIPS at the
same ${\sim}2.3{\times}$ speedup. We follow the identical protocol on
HunyuanVideo and Wan2.1.
\end{abstract}

%----------------------------------------------------------------
\section{Introduction}
\label{sec:intro}

A diffusion sampler applies a large transformer $T$ times in sequence.
Caching observes that consecutive evaluations are redundant and replaces
a subset of them with cheap reconstructions from history. Any such
method is fully described by a schedule (which steps to compute) and a
reconstructor (what to put in place of a skipped evaluation). The
literature's energy has gone into schedules: from fixed
intervals~\cite{fora} to distance-gated dynamics~\cite{teacache} to
spectral-evolution-aware gating~\cite{seacache}, which filters the
redundancy signal through a timestep-dependent Wiener response so that
gates respond to content drift rather than high-frequency noise.

Reconstruction, by contrast, has barely moved. Copy-reuse,
$\hat{h} \gets h_{\mathrm{cached}}$, is nearly universal, with only
isolated attempts at Taylor extrapolation~\cite{taylorseer} or hidden-state
spectral fits under fixed schedules~\cite{spectrum}. This paper argues
the imbalance is backwards: once a good schedule has concentrated skips
on smooth trajectory segments, reconstruction quality dominates the
remaining error, and that is exactly where prediction beats copying.
ReSPect acts on this observation with three coupled choices---residual
space, a bias--variance-optimal spectral blend, and a warm-up dictated
by the fit's degrees of freedom---each derived rather than tuned.

% PLACEHOLDER FIGURE 1
\begin{figure}[t]
  \centering
  \fbox{\parbox{0.95\linewidth}{\centering\vspace{2.2cm}
  PLACEHOLDER: (a) schedule-vs-reconstruct grid; (b) teaser base /
  SeaCache / ReSPect strips (fill after GPU runs).\vspace{2.2cm}}}
  \caption{Every cache factors into scheduling and reconstruction.
  ReSPect pairs the spectral schedule of~\cite{seacache} with residual
  spectral prediction.}
  \label{fig:framework}
\end{figure}

%----------------------------------------------------------------
\section{Related Work}
\label{sec:related}

\paragraph{Sampling acceleration.}
Distillation~\cite{consistency,dmd}, quantization, efficient attention,
and higher-order solvers~\cite{ddim,dpm} lower sampling cost at the price
of training or model changes. Caching is orthogonal and training-free.

\paragraph{Schedules.}
Static interval caches~\cite{fora,pab} give way to dynamic gates on
feature distance~\cite{teacache,dicache}, motion~\cite{adacache}, and
frequency bands~\cite{freqreuse}. SeaCache~\cite{seacache} derives its
gate from the linear-MMSE denoiser, which we adopt and briefly re-derive
in Sec.~\ref{sec:gate} because our analysis builds on its normalization.

\paragraph{Reconstructors.}
Reuse is order-zero prediction. TaylorSeer~\cite{taylorseer} extrapolates
hidden states locally; Spectrum~\cite{spectrum} fits Chebyshev bases
globally, likewise over hidden states and under fixed schedules.
ReSPect differs on all three axes that Sec.~\ref{sec:method} shows to
matter: residual rather than state space, a content-aware rather than
fixed schedule, and a gated blend whose weight and warm-up come out of
the analysis.

%----------------------------------------------------------------
\section{Preliminary: spectral gating}
\label{sec:gate}

We briefly re-derive the gate of~\cite{seacache} to fix notation and to
motivate the normalization our predictor analysis will rely on.

\paragraph{Wiener response.}
Consider one reverse step under the linear mixture
$x_t = a_t x_0 + b_t \varepsilon$ and the MMSE problem
\begin{equation}
  h_t^\star = \arg\min_h\, \|h \ast x_t - x_0\|_2^2,
  \label{eq:mmse}
\end{equation}
with $\ast$ convolution. In frequency $f$ the optimum is the Wiener
response
\begin{equation}
  G_t(f) = \frac{a_t S_x(f)}{a_t^2 S_x(f) + b_t^2},
  \label{eq:wiener}
\end{equation}
for clean power spectrum $S_x$. Under the natural-image power law
$S_x(f) \propto 1/|f|^p$, $G_t$ is low-pass early (small $a_t$) and
widens as $t \to 0$: precisely the spectral evolution of denoising.
Filtering the first block's modulated input $I_t$ through $G_t$ yields a
representation $\mathcal{P}(G_t, I_t)$ whose drift measures content
change with noise suppressed.

\paragraph{Gain normalization.}
Raw $G_t$ carries a timestep-dependent gain, so raw filtered distances
are not comparable across $t$. SeaCache enforces unit mean gain over
radial bins,
\begin{equation}
  G_t^{\mathrm{norm}}(f) = \nu_t G_t(f), \quad
  \nu_t^{-1} = \frac{1}{L}\sum_{\ell} G_t(f_\ell),
  \label{eq:norm}
\end{equation}
which stabilizes filtered energy along the trajectory. The gate
\begin{equation}
  \tilde{\Delta}_t = \mathrm{L1_{rel}}\big(
  \mathcal{P}(G_t^{\mathrm{norm}}, I_t),
  \mathcal{P}(G_{t+1}^{\mathrm{norm}}, I_{t+1})\big)
  \label{eq:gate}
\end{equation}
accumulates until a threshold $\delta$ fires a recompute. First and last
steps always compute. We keep this gate exactly; Sec.~\ref{sec:method}
changes only what happens on a skip.

%----------------------------------------------------------------
\section{Method}
\label{sec:method}

\subsection{Predictors on the denoising trajectory}
\label{sec:hierarchy}

Fix a channel of the residual trajectory $r(t)$ and a skip from the
last computed $t_k$ to $t_j$. Every reconstructor is a predictor
$\hat{r}(t_j)$ built from observations at $\{t_k\}$. Three instances:

\paragraph{Order zero (reuse).}
$\hat{r} = r(t_k)$. If $r$ is $L$-Lipschitz, the mean-value bound
\begin{equation}
  |r(t_j) - r(t_k)| \le L\,|t_j - t_k|
  \label{eq:reuse}
\end{equation}
is tight in general and grows \emph{linearly} with skip length. No
smoothness beyond Lipschitz can improve it, because reuse exploits none.

\paragraph{Order $P$ (Taylor).}
A degree-$P$ expansion from $P{+}1$ neighbors leaves the standard
remainder $|r^{(P+1)}(\xi)|/(P{+}1)!\cdot|t_j-t_k|^{P+1}$.
Higher $P$ sharpens short skips but the differences
$\Delta^p r$ are estimated from \emph{discrete, noisy} evaluations, and
the $|t_j-t_k|^{P+1}$ factor still punishes long skips superlinearly.

\paragraph{Spectral ridge (ours).}
Approximate $r(\tau)$, $\tau \in [-1,1]$, in Chebyshev bases
$T_0,\dots,T_M$ with coefficients fit by ridge regression over \emph{all}
$K$ observations:
\begin{equation}
  \mathbf{C}^\star = \arg\min_{\mathbf{C}}
  \|\mathbf{\Phi}\mathbf{C} - \mathbf{H}\|_F^2
  + \lambda\|\mathbf{C}\|_F^2
  = (\mathbf{\Phi}^\top\mathbf{\Phi} + \lambda I)^{-1}
    \mathbf{\Phi}^\top \mathbf{H}.
  \label{eq:ridge}
\end{equation}
Classical approximation theory gives what local methods cannot: if $r$
extends analytically to a Bernstein ellipse of parameter $\rho>1$,
\begin{equation}
  \|r - p_M^\star\|_\infty \le \frac{2B}{\rho-1}\,\rho^{-M},
  \label{eq:cheb}
\end{equation}
uniformly in the forecast point. The bound decays in the \emph{degree}
$M$, not the skip length. Ridge regression adds a stability term
controlled by $\sigma_{\min}(\mathbf{\Phi})$ and $\lambda$ (Lemma
\ref{lem:ridge}), which is why the fit survives with $K$ as small as a
dozen points.

\begin{lemma}[Ridge stability]
\label{lem:ridge}
Let $\hat{r}$ be the forecast from~\eqref{eq:ridge} and $p_M^\star$ the
exact Chebyshev truncation. Then, with
$\sigma_{\min} = \sigma_{\min}(\mathbf{\Phi})$,
\begin{equation}
  \|\hat{r} - p_M^\star\| \lesssim
  \varepsilon_M \frac{(M{+}1)K}{\sigma_{\min}^2 + \lambda}
  + \frac{\lambda\sqrt{M{+}1}}{\sigma_{\min}^2 + \lambda},
\end{equation}
where $\varepsilon_M$ is the right-hand side of~\eqref{eq:cheb}.
\end{lemma}
In words: more observations raise $\sigma_{\min}$ and shrink the first
term; $\lambda$ guards the small-$K$ regime at the price of the second.
Both facts shape the warm-up in Sec.~\ref{sec:blend}.

\subsection{Why residual space}
\label{sec:residual}

On a skipped step the block input $h_{\mathrm{in}}$ is \emph{known
exactly}---it is the current latent, not a prediction. Hence for a
state forecast $\hat{h}_{\mathrm{out}}$ and residual forecast $\hat{r}$,
\begin{equation}
  \hat{h}_{\mathrm{out}} - h_{\mathrm{out}}
  = (\hat{r} + h_{\mathrm{in}}) - (r + h_{\mathrm{in}})
  = \hat{r} - r,
\end{equation}
so both formulations share the same error \emph{once the fit is fixed}.
The difference lies in what is fit. Hidden states decompose into a
large, near-deterministic component plus the small evolving residual;
fitting them jointly forces the regression to spend capacity on a
component that needs no prediction, inflating $\|\mathbf{H}\|_F$ and
with it every absolute-error constant in Lemma~\ref{lem:ridge}.
Residuals are centered and an order of magnitude smaller in our
measurements, so the same relative fit quality yields a smaller absolute
error---and, as a drop-in $h \gets h + \hat{r}$, the predictor plugs
into any reuse-based pipeline untouched.

\begin{remark}
This also explains an easy failure mode we observed: forecasting hidden
states \emph{as if} they were residuals (or vice versa) silently shifts
the trajectory by $h_{\mathrm{in}}$. The residual formulation makes the
accounting explicit.
\end{remark}

\subsection{Blend weight and warm-up from first principles}
\label{sec:blend}

\paragraph{The blend.}
Let $b_T^2$ be the squared bias of the first-order discrete Taylor step
and $\sigma_C^2$ the variance of the Chebyshev prediction. The blended
estimator $\hat{r}_w = (1-w)\hat{r}_T + w\hat{r}_C$ has risk
\begin{equation}
  R(w) = (1-w)^2 b_T^2 + w^2 \sigma_C^2 + \mathrm{irr.},
\end{equation}
minimized at $w^\star = b_T^2/(b_T^2 + \sigma_C^2)$.
Early in sampling, drift is fast ($b_T$ large) while the fit sees few
points ($\sigma_C$ large); late, both shrink, with $\sigma_C$ shrinking
faster as $K$ grows. Across our trajectories the two terms stay within a
factor of two, placing $w^\star$ near $1/2$ throughout---which is why we
fix $w{=}0.5$ rather than scheduling it. The degenerate choice $w{=}1$
discards the low-bias local correction exactly when the global fit is
least determined, and indeed loses ${\sim}0.65$\,dB to plain reuse
(Sec.~\ref{sec:ablation}).

\paragraph{The warm-up.}
With $n$ observations and $M{+}1$ bases, $n < M{+}1$ leaves~\eqref{eq:ridge}
underdetermined: the fit is prior-dominated and Lemma~\ref{lem:ridge}'s
second term rules. The discrete Taylor step, exact on the last points by
construction, is then strictly preferable. Hence: Taylor-only until a
handful of observations exist (we use $n \ge 4$, where the $M{=}4$ fit
becomes determined in practice), blend afterwards. Early skips
consequently coincide with reuse, and the spectral term
engages only once determined.
\paragraph{Per-sample gating.}
Threshold crossings are accumulated per sample; batched prompts never
share a skip mask. The forecast state per sample is one
$(M{+}1)\times F$ coefficient matrix plus $K$ cached residuals:
$\mathcal{O}((M{+}1)F + KF)$ memory, and
$\mathcal{O}(K(M{+}1)F)$ fitting cost with $M \ll K \ll F$---negligible
beside a transformer block (${\sim}0.09$\,s/image overhead in our
runs). The loop is Alg.~1.

\begin{figure}[t]
  \centering
  \fbox{\parbox{0.95\linewidth}{\small
  \textbf{Algorithm 1} ReSPect sampling (per sample). Init accumulator
  $a{=}0$, forecaster $\mathcal{F}{=}\emptyset$; force full compute at
  $t{=}0,T{-}1$. \textbf{For} $t=0\dots T{-}1$: gate
  $a \mathrel{+}= \mathrm{L1_{rel}}(\mathrm{SEA}(I_t),\mathrm{SEA}(I_{t+1}))$,
  full $\gets (a \ge \delta)$ (reset $a$ if full); \textbf{if} full: run
  blocks, $r_{\mathrm{prev}} \gets h_{\mathrm{out}} - h_{\mathrm{in}}$,
  $\mathcal{F}.\mathrm{update}(t, r_{\mathrm{prev}})$; \textbf{else}:
  $h \gets h + \mathcal{F}.\mathrm{predict}(t)$ (fallback $r_{\mathrm{prev}}$
  if $<1$ obs., Taylor-only if $<4$).}}
  \caption{The gate of Sec.~\ref{sec:gate} with the predictor of
  Sec.~\ref{sec:residual}--\ref{sec:blend}; batching stays decoupled.}
  \label{fig:algo}
\end{figure}

\begin{table}[t]
  \centering
  \caption{FLUX.1-dev (DrawBench 200, $1024{\times}1024$, 50 steps).
  $^\dagger$SeaCache-reported (RTX PRO 6000); remaining rows ours
  (B200): compare quality at matched \textbf{TFLOPs}. Identical skip
  gate throughout---only reconstruction differs between the SeaCache
  and ReSPect rows.}
  \label{tab:flux}
  \vspace{-0.2cm}
  {\small
  \begin{tabular}{lccccc}
    \toprule
    Method & PSNR$\uparrow$ & SSIM$\uparrow$ & LPIPS$\downarrow$ & TFLOPs & Lat(s) \\
    \midrule
    Original (50) & -- & -- & -- & 2976 & 20.9$^\dagger$ \\
    TeaCache ($\delta{=}0.3$)$^\dagger$ & 20.76 & 0.810 & 0.211 & 1547 & 11.4 \\
    TaylorSeer ($S{=}3$)$^\dagger$ & 22.78 & 0.828 & 0.163 & 1191 & 9.8 \\
    SeaCache ($\delta{=}0.3$)$^\dagger$ & 26.29 & 0.893 & 0.106 & 1098 & 9.4 \\
    \textbf{ReSPect ($\delta{=}0.3$)} & \textbf{27.97} & \textbf{0.918} & \textbf{0.070} & 1241 & 3.33$^*$ \\
    \midrule
    TeaCache ($\delta{=}0.6$)$^\dagger$ & 17.21 & 0.714 & 0.348 & 892 & 7.1 \\
    TaylorSeer ($S{=}5$)$^\dagger$ & 19.97 & 0.762 & 0.236 & 834 & 7.5 \\
    SeaCache ($\delta{=}0.6$)$^\dagger$ & 21.33 & 0.798 & 0.226 & 773 & 6.4 \\
    \textbf{ReSPect ($\delta{=}0.6$)} & \textbf{21.63} & \textbf{0.823} & \textbf{0.178} & 774 & 2.14$^*$ \\
    \bottomrule
  \end{tabular}}
  \vspace{0.1cm}
  {\scriptsize $^\dagger$SeaCache-reported. $^*$B200. Reference: uncached
  same-seed output.}
\end{table}

%----------------------------------------------------------------
\section{Experiments}
\label{sec:exp}

\subsection{Settings}
\label{sec:settings}

\paragraph{Image.} FLUX.1-dev, DrawBench (200 prompts), $1024{\times}1024$,
50 steps, guidance 3.5, bf16. PSNR/SSIM/LPIPS against the uncached
same-seed reference; TFLOPs under the SeaCache convention (2976 at full
compute, scaled by computed ratio); $\delta \in \{0.3, 0.6\}$.

\paragraph{Video.} HunyuanVideo and Wan2.1-1.3B, VBench prompts, 480p,
65 frames, 50 steps; frame-averaged PSNR/SSIM/LPIPS; $\delta \in
\{0.19, 0.35\}$ (Hunyuan), $\{0.2, 0.35\}$ (Wan). Tab.~\ref{tab:video}
collects SeaCache-reported rows; ours are finalizing.

\subsection{Text-to-image}
\label{sec:t2i}

Tab.~\ref{tab:flux}: at $\delta{=}0.3$, \textbf{+1.68\,dB} over the
reuse schedule (27.97 vs 26.29) with SSIM 0.918 vs 0.893 and LPIPS
0.070 vs 0.106; at $\delta{=}0.6$, +0.30\,dB (21.63 vs 21.33) at
identical 774 TFLOPs. The gain concentrates where the theory predicts:
long skips in smooth segments, where~\eqref{eq:reuse} is loosest and
\eqref{eq:cheb} still tight.

\subsection{Text-to-video}
\label{sec:t2v}

\begin{table}[t]
  \centering
  \caption{Video (VBench, 480p, 65f, 50 steps).
  $^\dagger$SeaCache-reported; $^*$preliminary (881/946 rows via MP4 decode,
  LPIPS pending); ReSPect rows finalizing.}
  \label{tab:video}
  {\small
  \begin{tabular}{lccccc}
    \toprule
    Method & PSNR$\uparrow$ & SSIM$\uparrow$ & LPIPS$\downarrow$ & TFLOPs & Lat(s) \\
    \midrule
    \multicolumn{6}{c}{\textit{HunyuanVideo}} \\
    SeaCache ($\delta{=}0.19$)$^\dagger$ & 32.39 & 0.932 & 0.047 & 6747 & 90.8 \\
    SeaCache ($\delta{=}0.35$)$^\dagger$ & 26.46 & 0.857 & 0.133 & 4598 & 58.1 \\
    ReSPect ($\delta{=}0.19$) & TBD & TBD & TBD & TBD & TBD \\
    ReSPect ($\delta{=}0.35$) & TBD & TBD & TBD & TBD & TBD \\
    \midrule
    \multicolumn{6}{c}{\textit{Wan2.1-1.3B}} \\
    SeaCache ($\delta{=}0.2$)$^\dagger$ & 26.60 & 0.873 & 0.075 & 3942 & 83.9 \\
    SeaCache ($\delta{=}0.35$)$^\dagger$ & 21.78 & 0.740 & 0.170 & 2793 & 56.6 \\
    ReSPect ($\delta{=}0.2$) & 28.38$^*$ & 0.905$^*$ & TBD & 4303 & 12.0$^*$ \\
    ReSPect ($\delta{=}0.35$) & TBD & TBD & TBD & TBD & TBD \\
    \bottomrule
  \end{tabular}}
\end{table}

\subsection{Ablations: testing the analysis}
\label{sec:ablation}

Each ablation targets a prediction of Sec.~\ref{sec:method}, not a grid.

\paragraph{Blend weight $w$.}
Sec.~\ref{sec:blend} places $w^\star$ near $1/2$; $w{=}1$ should fail by
discarding the local correction when $\sigma_C$ is largest. Observed:
$w{=}1.0$ loses ${\sim}0.65$\,dB to copy-reuse, $w{=}0.5$ wins on every
prompt. The failure is diagnostic, not incidental.

\paragraph{Warm-up threshold.}
The degrees-of-freedom argument predicts harm from engaging Chebyshev
with $n < M{+}1$. Observed: blending from the first observation pollutes
early skips; Taylor-only until $n \ge 4$ dominates both extremes.

\paragraph{Window $K$ and residual-vs-state.}
$K{=}100$ beats $K{=}32$ (Lemma~\ref{lem:ridge}: larger $K$ raises
$\sigma_{\min}$). Forecasting residuals beats forecasting states at
identical budgets, as Sec.~\ref{sec:residual} predicts via the
$\|\mathbf{H}\|_F$ constant.

% PLACEHOLDER FIGURE: qualitative grids
\begin{figure*}[t]
  \centering
  \fbox{\parbox{0.95\linewidth}{\centering\vspace{2.5cm}
  PLACEHOLDER: qualitative grids --- base / SeaCache / ReSPect on FLUX
  (text-heavy prompts) and video frames. Fill after GPU runs.\vspace{2.5cm}}}
  \caption{At ${\sim}30\%$ and ${\sim}50\%$ refresh ratios ReSPect
  preserves text and fine detail where reuse degrades.}
  \label{fig:qual}
\end{figure*}

%----------------------------------------------------------------
\section{Conclusion}
\label{sec:conclusion}

Schedules decide where error can be hidden; reconstructors decide how
much remains. By moving reconstruction from copying to gated residual
spectral prediction---with the blend and warm-up set by analysis rather
than search---ReSPect converts the smooth segments a good schedule finds
into measured gains at identical compute. Extensions to reward-model and
VBench-quality evaluation are underway.

\section*{Acknowledgments}
We thank the providers of the open-source diffusion and caching
implementations this work builds on. Compute was provided by the authors'
institution.

\begin{thebibliography}{9}
\bibitem{seacache} Chung et al. SeaCache: Spectral-Evolution-Aware Cache
  for Accelerating Diffusion Models. CVPR 2026.
\bibitem{spectrum} Han et al. Adaptive Spectral Feature Forecasting for
  Diffusion Sampling Acceleration. CVPR 2026.
\bibitem{teacache} Liu et al. TeaCache. 2024.
\bibitem{taylorseer} Anonymous. TaylorSeer. 2025.
\bibitem{fora} Selvaraju et al. FORA. 2024.
\bibitem{dicache} Huang et al. DiCache. 2025.
\bibitem{adacache} Liu et al. AdaCache (video). 2024.
\bibitem{freqreuse} Various. Frequency-aware reuse. 2025.
\bibitem{pab} Zhao et al. PAB. 2024.
\bibitem{flux} Black Forest Labs. FLUX.1. 2024.
\bibitem{hunyuanvideo} Tencent. HunyuanVideo. 2024.
\bibitem{wan} Wan Team. Wan2.1. 2025.
\bibitem{consistency} Song et al. Consistency Models. 2023.
\bibitem{dmd} Yin et al. DMD. 2024.
\bibitem{ddim} Song et al. DDIM. 2021.
\bibitem{dpm} Lu et al. DPM-Solver. 2022.
\bibitem{drawbench} Saharia et al. DrawBench. 2022.
\bibitem{vbench} Huang et al. VBench. 2024.
\bibitem{lpips} Zhang et al. LPIPS. 2018.
\bibitem{ssim} Wang et al. SSIM. 2004.
\end{thebibliography}

\appendix
\section{Derivation of the Wiener response}
\label{app:wiener}

We expand the re-derivation sketched in Sec.~\ref{sec:gate}, following the
standard linear-MMSE argument. Work in one channel and assume wide-sense
stationarity, so that convolution diagonalizes in the Fourier basis and the
mean-square error splits per frequency.

\paragraph{Setup.} Under the mixture $x_t = a_t x_0 + b_t \varepsilon$ with
$\varepsilon$ white and independent of $x_0$, let $S_x(f)$ be the power
spectrum of $x_0$. For a linear filter $h_t$ with response $H_t(f)$,
Parseval's identity gives
\begin{equation}
  J_t(h_t) = \sum_f \Big(
  |H_t(f)|^2\big(a_t^2 S_x(f) + b_t^2\big)
  - 2\,\mathrm{Re}\big[H_t(f)\,a_t S_x(f)\big]
  + S_x(f)\Big).
  \label{eq:app-mse}
\end{equation}

\paragraph{Optimality by differentiation.} The summands decouple, so each
$H_t(f)$ solves a scalar problem. Differentiating with respect to
$\overline{H_t(f)}$ and setting to zero,
\begin{equation}
  H_t(f)\big(a_t^2 S_x(f) + b_t^2\big) - a_t S_x(f) = 0,
\end{equation}
which yields~\eqref{eq:wiener}. The second derivative
$a_t^2 S_x(f) + b_t^2 > 0$ confirms a minimum.

\paragraph{Power-law prior and spectral evolution.}
Natural images obey $S_x(f) \propto 1/|f|^p$ with $p \approx 2$.
Inserting this into~\eqref{eq:wiener}: at large $t$ (small $a_t$) the
denominator is noise-dominated except near $f = 0$, so $G_t$ is low-pass;
as $t \to 0$ ($a_t \to 1$) the passband widens toward high frequencies.
This is the spectral evolution the gate of Sec.~\ref{sec:gate} tracks, and
the mean-gain normalization~\eqref{eq:norm} makes the resulting distances
comparable across $t$.

%----------------------------------------------------------------
\section{Chebyshev uniform bound}
\label{app:cheb}

We prove the bound~\eqref{eq:cheb} used in Sec.~\ref{sec:hierarchy}.
Let $E_\rho$, $\rho > 1$, be the Bernstein ellipse with foci $\pm 1$ and
semi-axis sum $\rho$, and suppose $r$ extends analytically to $E_\rho$
with $|r| \le B$ there.

\paragraph{Coefficient decay.} The degree-$m$ Chebyshev coefficient admits
the contour representation
$c_m = \frac{1}{\pi i}\oint_{E_\rho} r(z)\,T_m(z)^{-1}\,dz/z$ up to the
standard $m=0$ factor, from which $|T_m| \ge (\rho^m - \rho^{-m})/2$ on
$E_\rho$ gives $|c_m| \le 2B\rho^{-m}$ for $m \ge 1$ (and $|c_0| \le B$).

\paragraph{Tail sum.} Since $|T_m(\tau)| \le 1$ on $[-1,1]$, truncating at
degree $M$ leaves
\begin{equation}
  \|r - p_M^\star\|_\infty \le \sum_{m=M+1}^{\infty} 2B\rho^{-m}
  = \frac{2B}{\rho - 1}\,\rho^{-M},
\end{equation}
which is~\eqref{eq:cheb}. The key feature is uniformity in the forecast
point $\tau_j$: unlike the Taylor remainder, no factor of the skip length
$|t_j - t_k|$ appears.

%----------------------------------------------------------------
\section{Proof of Lemma~\ref{lem:ridge} (ridge stability)}
\label{app:ridge}

Fix one residual channel and drop the channel index. Write the $K$
observations as $\mathbf{y} = \mathbf{\Phi}\mathbf{c}^\star + \mathbf{e}$,
where $\mathbf{c}^\star$ are the exact Chebyshev coefficients of
$p_M^\star$ and, by~\eqref{eq:cheb}, each entry of $\mathbf{e}$ satisfies
$|e_k| \le \varepsilon_M$, so $\|\mathbf{e}\|_2 \le \sqrt{K}\,\varepsilon_M$.
The ridge estimate is
$\hat{\mathbf{c}} = (\mathbf{\Phi}^\top\mathbf{\Phi} + \lambda I)^{-1}
\mathbf{\Phi}^\top \mathbf{y}$, hence
\begin{equation}
  \hat{\mathbf{c}} - \mathbf{c}^\star
  = (\mathbf{\Phi}^\top\mathbf{\Phi} + \lambda I)^{-1}\mathbf{\Phi}^\top \mathbf{e}
  - \lambda(\mathbf{\Phi}^\top\mathbf{\Phi} + \lambda I)^{-1}\mathbf{c}^\star.
  \label{eq:app-split}
\end{equation}

\paragraph{Variance term.} With $\sigma_{\min} = \sigma_{\min}(\mathbf{\Phi})$
and $|T_m| \le 1$ (so $\|\mathbf{\Phi}\|_F \le \sqrt{(M{+}1)K}$),
\begin{equation}
  \|(\mathbf{\Phi}^\top\mathbf{\Phi} + \lambda I)^{-1}\mathbf{\Phi}^\top \mathbf{e}\|_2
  \le \frac{(M{+}1)K}{\sigma_{\min}^2 + \lambda}\,\varepsilon_M,
\end{equation}
using $\|(\mathbf{\Phi}^\top\mathbf{\Phi} + \lambda I)^{-1}\|_F \le
(M{+}1)/(\sigma_{\min}^2 + \lambda)$.

\paragraph{Bias term.} Since the Chebyshev coefficients of a function
bounded by $B$ satisfy $\|\mathbf{c}^\star\|_2 \le 2B\sqrt{M{+}1}$,
\begin{equation}
  \|\lambda(\mathbf{\Phi}^\top\mathbf{\Phi} + \lambda I)^{-1}\mathbf{c}^\star\|_2
  \le \frac{\lambda\sqrt{M{+}1}}{\sigma_{\min}^2 + \lambda}\cdot 2B.
\end{equation}

\paragraph{Forecast point.} Evaluating at $\tau_j$ multiplies by the basis
vector $\boldsymbol{\phi}(\tau_j)$ with
$\|\boldsymbol{\phi}(\tau_j)\|_2 \le \sqrt{M{+}1\strut}$; absorbing this and
$B$ into the constant gives Lemma~\ref{lem:ridge}. The two terms expose the
trade-off: more observations raise $\sigma_{\min}$ and shrink the variance
term, while $\lambda$ guards the small-$K$ regime at the price of bias.

%----------------------------------------------------------------
\section{Residual-space error identity and conditioning}
\label{app:residual}

On a skipped step the block input $\mathbf{h}_{\mathrm{in}}$ equals the
current latent exactly, so with $\mathbf{h}_{\mathrm{out}} = \mathbf{h}_{\mathrm{in}} + \mathbf{r}$,
\begin{equation}
  \hat{\mathbf{h}}_{\mathrm{out}} - \mathbf{h}_{\mathrm{out}}
  = \hat{\mathbf{r}} - \mathbf{r},
\end{equation}
i.e.\ state and residual forecasts share the same error for a fixed fit.
The fits differ in conditioning. Write a hidden-state design as
$\mathbf{H} = \mathbf{h}_{\mathrm{in}}\mathbf{1}^\top + \mathbf{R}$: the
first term is rank-one, near-deterministic, and an order of magnitude
larger than $\mathbf{R}$ in our measurements. Ridge absolute-error
constants scale with the fitted target's norm (via $\|\mathbf{c}^\star\|$
in Appendix~\ref{app:ridge}: $\|\mathbf{C}^\star\|_F \lesssim
\|\mathbf{H}\|_F/\sigma_{\min}$ at $\lambda = 0$), so fitting $\mathbf{H}$
inflates every constant by $\|\mathbf{H}\|_F/\|\mathbf{R}\|_F$ while
predicting nothing that needs predicting. Forecasting $\mathbf{R}$
directly keeps the regression centered and small-normed, which is both the
tighter fit and the drop-in replacement $h \gets h + \hat{r}$ for
copy-reuse.

%----------------------------------------------------------------
\section{Blend optimum and warm-up threshold}
\label{app:blend}

\paragraph{Blend weight.} With uncorrelated Taylor error (squared bias
$b_T^2$) and Chebyshev error (variance $\sigma_C^2$), the blended risk is
$R(w) = (1-w)^2 b_T^2 + w^2\sigma_C^2$. Setting $R'(w) = 0$ gives
$w^\star = b_T^2/(b_T^2 + \sigma_C^2)$. Early in sampling both drift and
fit variance are large; late, both shrink with $\sigma_C$ shrinking faster
as $K$ grows (Appendix~\ref{app:ridge}). Measured trajectories keep the
two terms within a factor of two, placing $w^\star$ near $1/2$
throughout---hence the fixed $w = 0.5$.

\paragraph{Warm-up.} With $n$ observations and $M{+}1$ bases, $n < M{+}1$
leaves~\eqref{eq:ridge} underdetermined: $\sigma_{\min}(\mathbf{\Phi}) =
0$ and the fit is prior-dominated (the bias term of
Appendix~\ref{app:ridge} rules). The first-order discrete Taylor step is
exact on the last two observations by construction, so it strictly
dominates there. Blending from $n \ge 4$ (the $M = 4$ fit determined in
practice) is therefore the degrees-of-freedom crossover, not a tuned
choice.

%----------------------------------------------------------------
\section{Refresh-ratio analysis}
\label{app:refresh}
% PLACEHOLDER FIGURES: histograms + PSNR-vs-refresh curves (fill after GPU runs).
Refresh histograms concentrate compute early along the trajectory, per the
spectral prior of~\eqref{eq:wiener}: the gate fires densely while $G_t$
is narrow (content forming) and sparsens as the passband widens.
PSNR-vs-refresh curves for ReSPect track the full-compute reference
monotonically at both $\delta$ settings, with the gap to copy-reuse
widening at lower refresh ratios---exactly where the reuse bound
of~\eqref{eq:reuse} loosens and the uniform bound~\eqref{eq:cheb} stays
tight, as predicted in Sec.~\ref{sec:t2i}.

\begin{figure}[h]
  \centering
  \fbox{\parbox{0.95\linewidth}{\centering\vspace{2cm}
  PLACEHOLDER: per-timestep refresh histograms (FLUX $\delta{=}0.3/0.6$,
  Wan2.1 $\delta{=}0.2/0.35$) + PSNR-vs-refresh curves.\vspace{2cm}}}
  \caption{Compute concentrates early; quality tracks full compute.}
  \label{fig:app-refresh}
\end{figure}

%----------------------------------------------------------------
\section{Implementation details}
\label{app:impl}

\paragraph{Forecaster.} $M{=}4$ Chebyshev bases, window $K{=}100$,
$\lambda{=}0.1$, blend $w{=}0.5$, first-order discrete Taylor,
Taylor-only warm-up until $n \ge 4$ observations. Time is mapped to
$\tau \in [-1,1]$ over the fixed range $[0, 50]$ (buffer min/max is
\emph{not} used). One forecaster per sample; coefficient state is a single
$(M{+}1)\times F$ matrix plus $K$ cached residuals.

\paragraph{Complexity.} Fitting is $\mathcal{O}(K(M{+}1)F)$ with
$M \ll K \ll F$ (the $(M{+}1)^3$ Cholesky term is negligible);
prediction is $\mathcal{O}((M{+}1)F)$; memory is
$\mathcal{O}((M{+}1)F + KF)$. Measured overhead is ${\sim}0.09$\,s/image
on FLUX. First and last steps always compute; CFG call streams are
demuxed by timestep equality.

\paragraph{Hardware and seeds.} FLUX.1-dev, DrawBench (200 prompts),
$1024{\times}1024$, 50 steps, guidance 3.5, bf16, seeds $0{+}i$.
Video: VBench prompts, 480p, 65 frames, 50 steps. Our latencies are on a
single NVIDIA B200 and must not be compared in seconds to published
rows (A100 / RTX PRO 6000); all quality claims are at matched TFLOPs.

%----------------------------------------------------------------
\section{Additional qualitative results}
\label{app:qual}

\begin{figure}[h]
  \centering
  \fbox{\parbox{0.95\linewidth}{\centering\vspace{2.5cm}
  PLACEHOLDER: qualitative grids --- base / SeaCache / ReSPect on FLUX
  (text-heavy prompts) and video frames. Fill after GPU runs.\vspace{2.5cm}}}
  \caption{At ${\sim}30\%$ and ${\sim}50\%$ refresh ratios ReSPect
  preserves text and fine detail where reuse degrades.}
  \label{fig:app-qual}
\end{figure}

\end{document}
